
\subsubsection{2.1 Proof Checker Feedback}

\textbf{BFS-Prover} \citep{Xinetal2025} combines Best-First Tree Search with DPO training using Lean4 compiler errors as rejected examples. At each proof state, the Lean4 compiler returns structured error messages for failed tactics; these become the rejected examples in DPO pairs. The system achieves 72.95\% on miniF2F. DPO hyperparameters include β=10, learning rate decay from 5e-6 to 5e-7, and a 1-epoch training schedule. BFS-Prover does not compare its feedback condition against semantically ungrounded alternatives, so whether the semantic content of compiler errors (versus simply having same-state failed tactics as negatives) drives the improvement is unknown.

\textbf{STP} [Dong and Ma, 2025] achieves 65\% on miniF2F using self-play with formal verifier feedback at the episode level. Like BFS-Prover, STP uses a single feedback type without cross-condition ablation.

\textbf{Neural Theorem Proving} \citep{Raoetal2025} applies SFT+RL with Isabelle verifier feedback, achieving state-of-the-art results on Isabelle benchmarks, again without testing the mechanistic question.

\subsubsection{2.2 SMT Solver Feedback}

\textbf{Proof of Thought} \citep{Gangulyetal2024} uses Z3 SMT counterexamples via a JSON DSL bridge in an inference-time repair loop (no DPO training), reducing compilation errors from 14.6\% to 0\%. The Z3 counterexample provides episode-level grounded feedback, not step-local guidance.

\textbf{Step-Wise Formal Verification} [Zhou and Zhang, 2025] uses SMT solvers and computer algebra systems as external oracles for LLM math verification at inference time.

\subsubsection{2.3 Static Analysis Feedback}

\textbf{PropertyGPT} \citep{Liuetal2024} uses PSL compiler feedback to guide LLM generation of smart contract formal properties, improving compilation success from 63\% to 87\% and achieving 80\% recall on Certora-audited projects. This is the strongest available inference-time evidence consistent with the oracle framing. Structured feedback with explicit semantic content outperforms no-feedback baselines; the mechanism is not isolated.

\subsubsection{2.4 DPO for Formal Reasoning}

\textbf{Step-DPO} \citep{Laietal2024} applies step-level preference optimization to informal mathematical reasoning, treating individual reasoning steps as preference units. This is structurally analogous to tactic-level DPO in formal verification but operates without a formal verifier.

\subsubsection{2.5 Benchmark Landscape}

\textbf{miniF2F} \citep{Zhengetal2021} provides 488 theorem proving problems (AMC/AIME/IMO and high-school/undergraduate mathematics) in Lean4 format. \textbf{Vericoding} \citep{Bursucetal2025} provides 12,504 formally verified code specifications across Lean4/Verus/Dafny formalisms. Neither benchmark includes feedback-loop evaluation metrics such as locality score or feedback signal quality.

\subsubsection{2.6 Positioning}

No prior work introduces the oracle/regularizer framing as a testable distinction, proposes a mechanistic metric (locality score) to directly measure oracle function, or designs a controlled cross-condition experiment within a fixed BFS framework to separate semantic from structural contributions of formal feedback.
